% GENERATED FILE. Edit canonical lessons, then run npm run generate:books.
% canonical-source-hash: fnv1a64:89a92c28b8e2f03d
% canonical-lessons: AR-C468-nabat, AR-C468-shujayra, AR-C468-dhabala, AR-C468-azhara, AR-C468-athmara

\chapter{Plants}
\label{ch:ar-plants}
\hlchaptermodality{468}

\begin{chapteropening}
\textbf{By the end of this chapter:} \emph{I can say plant life, a shrub, to wilt, to bloom and to bear fruit.}

Complete the last lesson of chapter 468: I can say plant life, a shrub, to wilt, to bloom and to bear fruit.
\end{chapteropening}

\section[\texorpdfstring{\ar{نبات} (nabāt)}{نبات (nabāt)}]{\ar{نبات} (nabāt) — plant life; a plant}

\label{lesson:AR-C468-nabat}



\begin{quote}
Before the new one: say the Arabic for to die out, then the Arabic for to squander.
\end{quote}

\subsection*{You'll want to know: \ar{نبات}}
\textbf{\ar{نبات}} — \emph{nabāt} — \textquotedblleft{}plant life; a plant\textquotedblright{}.

\begin{quote}
\ar{رأينا} \ar{نباتا} \ar{غريبا} \ar{بجانب} \ar{الشلال}.
\end{quote}

\begin{quote}
\emph{raʾaynā nabātan gharīban bi-jānib ash-shallāl.}
\end{quote}

\begin{quote}
We saw a strange plant beside the waterfall.
\end{quote}

\begin{cousinweb}
\textbf{\ar{نبات}} (\emph{nabāt}) is built on the native Arabic root \emph{n-b-t}, \textquotedblleft{}to grow, to sprout\textquotedblright{} (Lane).
\end{cousinweb}

\begin{grammarlens}[title={What you've built}]
One more word for getting around.
\end{grammarlens}

\subsection*{Your turn}
\begin{itemize}
\raggedright
  \item \emph{Say it:} \emph{nabāt}
  \item \emph{Say it:} \emph{nabāt}, once more
  \item \emph{Say it:} \emph{raʾaynā nabātan gharīban bi-jānib ash-shallāl.} — We saw a strange plant beside the waterfall.
  \item \emph{Recall:} say the Arabic for to die out, then the Arabic for to squander, then say \emph{nabāt} again
\end{itemize}

\begin{tcolorbox}[breakable,colback=teal!4,colframe=teal!35!black,title={Before you move on}]
What does \ar{نبات} mean? (\textquotedblleft{}Plant life; a plant\textquotedblright{}.) Say it once more.
\end{tcolorbox}
\section[\texorpdfstring{\ar{شجيرة} (shujayra)}{شجيرة (shujayra)}]{\ar{شجيرة} (shujayra) — a shrub; a bush}

\label{lesson:AR-C468-shujayra}



\begin{quote}
Before the new one: say the Arabic for to squander, then the Arabic for plant life.
\end{quote}

\subsection*{You'll want to know: \ar{شجيرة}}
\textbf{\ar{شجيرة}} — \emph{shujayra} — \textquotedblleft{}a shrub; a bush\textquotedblright{}.

\begin{quote}
\ar{زرع} \ar{جدي} \ar{شجيرة} \ar{ورد} \ar{في} \ar{الحديقة}.
\end{quote}

\begin{quote}
\emph{zaraʿa jaddī shujayrat ward fī al-ḥadīqa.}
\end{quote}

\begin{quote}
My grandfather planted a rose bush in the garden.
\end{quote}

The pattern \emph{fuʿayla} makes a small one: \emph{shujayra} is a little \emph{shajara}, \textquotedblleft{}a tree\textquotedblright{}.

\begin{cousinweb}
\textbf{\ar{شجيرة}} (\emph{shujayra}) is built on the native Arabic root \emph{sh-j-r}, \textquotedblleft{}trees\textquotedblright{} (Lane).
\end{cousinweb}

\begin{grammarlens}[title={What you've built}]
One more word for getting around.
\end{grammarlens}

\subsection*{Your turn}
\begin{itemize}
\raggedright
  \item \emph{Say it:} \emph{shujayra}
  \item \emph{Say it:} \emph{shujayra}, once more
  \item \emph{Say it:} \emph{zaraʿa jaddī shujayrat ward fī al-ḥadīqa.} — My grandfather planted a rose bush in the garden.
  \item \emph{Recall:} say the Arabic for to squander, then the Arabic for plant life, then say \emph{shujayra} again
\end{itemize}

\begin{tcolorbox}[breakable,colback=teal!4,colframe=teal!35!black,title={Before you move on}]
What does \ar{شجيرة} mean? (\textquotedblleft{}A shrub; a bush\textquotedblright{}.) Say it once more.
\end{tcolorbox}
\section[\texorpdfstring{\ar{ذبل} (dhabala)}{ذبل (dhabala)}]{\ar{ذبل} (dhabala) — to wilt; to wither}

\label{lesson:AR-C468-dhabala}



\begin{quote}
Before the new one: say the Arabic for plant life, then the Arabic for a shrub.
\end{quote}

\subsection*{You'll want to know: \ar{ذبل}}
\textbf{\ar{ذبل}} — \emph{dhabala} — \textquotedblleft{}to wilt; to wither\textquotedblright{}.

\begin{quote}
\ar{ذبلت} \ar{الزهرة} \ar{في} \ar{الشمس}.
\end{quote}

\begin{quote}
\emph{dhabalat az-zahra fī ash-shams.}
\end{quote}

\begin{quote}
The flower wilted in the sun.
\end{quote}

Present: \ar{يذبل}, \emph{yadhbul}.

\begin{cousinweb}
\textbf{\ar{ذبل}} (\emph{dhabala}) is usually traced to the native Arabic root \emph{dh-b-l}, \textquotedblleft{}to wither\textquotedblright{}. It is the plain verb.
\end{cousinweb}

\begin{grammarlens}[title={What you've built}]
One more word for getting around.
\end{grammarlens}

\subsection*{Your turn}
\begin{itemize}
\raggedright
  \item \emph{Say it:} \emph{dhabala}
  \item \emph{Say it:} \emph{dhabala}, once more
  \item \emph{Say it:} \emph{dhabalat az-zahra fī ash-shams.} — The flower wilted in the sun.
  \item \emph{Recall:} say the Arabic for plant life, then the Arabic for a shrub, then say \emph{dhabala} again
\end{itemize}

\begin{tcolorbox}[breakable,colback=teal!4,colframe=teal!35!black,title={Before you move on}]
What does \ar{ذبل} mean? (\textquotedblleft{}To wilt; to wither\textquotedblright{}.) Say it once more.
\end{tcolorbox}
\section[\texorpdfstring{\ar{أزهر} (azhara)}{أزهر (azhara)}]{\ar{أزهر} (azhara) — to bloom; to flower}

\label{lesson:AR-C468-azhara}



\begin{quote}
Before the new one: say the Arabic for a shrub, then the Arabic for to wilt.
\end{quote}

\subsection*{You'll want to know: \ar{أزهر}}
\textbf{\ar{أزهر}} — \emph{azhara} — \textquotedblleft{}to bloom; to flower\textquotedblright{}.

\begin{quote}
\ar{أزهرت} \ar{الأشجار} \ar{في} \ar{الربيع}.
\end{quote}

\begin{quote}
\emph{azharat al-ashjār fī ar-rabīʿ.}
\end{quote}

\begin{quote}
The trees blossomed in spring.
\end{quote}

Present: \ar{يزهر}, \emph{yuzhir}.

\begin{cousinweb}
\textbf{\ar{أزهر}} (\emph{azhara}) is built on the native Arabic root \emph{z-h-r}, \textquotedblleft{}flowers; to shine\textquotedblright{} (Lane). It is Form IV, \textquotedblleft{}to put out flowers\textquotedblright{}; \emph{zahra}, \textquotedblleft{}a flower\textquotedblright{}, shares the root.
\end{cousinweb}

\begin{grammarlens}[title={What you've built}]
One more word for getting around.
\end{grammarlens}

\subsection*{Your turn}
\begin{itemize}
\raggedright
  \item \emph{Say it:} \emph{azhara}
  \item \emph{Say it:} \emph{azhara}, once more
  \item \emph{Say it:} \emph{azharat al-ashjār fī ar-rabīʿ.} — The trees blossomed in spring.
  \item \emph{Recall:} say the Arabic for a shrub, then the Arabic for to wilt, then say \emph{azhara} again
\end{itemize}

\begin{tcolorbox}[breakable,colback=teal!4,colframe=teal!35!black,title={Before you move on}]
What does \ar{أزهر} mean? (\textquotedblleft{}To bloom; to flower\textquotedblright{}.) Say it once more.
\end{tcolorbox}
\section[\texorpdfstring{\ar{أثمر} (athmara)}{أثمر (athmara)}]{\ar{أثمر} (athmara) — to bear fruit}

\label{lesson:AR-C468-athmara}



\begin{quote}
Before the new one: say the Arabic for to wilt, then the Arabic for to bloom.
\end{quote}

\subsection*{You'll want to know: \ar{أثمر}}
\textbf{\ar{أثمر}} — \emph{athmara} — \textquotedblleft{}to bear fruit\textquotedblright{}.

\begin{quote}
\ar{أثمرت} \ar{شجرة} \ar{التين} \ar{في} \ar{الصيف}.
\end{quote}

\begin{quote}
\emph{athmarat shajarat at-tīn fī aṣ-ṣayf.}
\end{quote}

\begin{quote}
The fig tree bore fruit in summer.
\end{quote}

Present: \ar{يثمر}, \emph{yuthmir}.

Work can bear fruit too: the sense carries over as in English.

\begin{cousinweb}
\textbf{\ar{أثمر}} (\emph{athmara}) is built on the native Arabic root \emph{th-m-r}, \textquotedblleft{}fruit\textquotedblright{} (Lane). It is Form IV, \textquotedblleft{}to put out fruit\textquotedblright{}.
\end{cousinweb}

\begin{grammarlens}[title={What you've built}]
One more word for getting around.
\end{grammarlens}

\subsection*{Your turn}
\begin{itemize}
\raggedright
  \item \emph{Say it:} \emph{athmara}
  \item \emph{Say it:} \emph{athmara}, once more
  \item \emph{Say it:} \emph{athmarat shajarat at-tīn fī aṣ-ṣayf.} — The fig tree bore fruit in summer.
  \item \emph{Recall:} say the Arabic for to wilt, then the Arabic for to bloom, then say \emph{athmara} again
\end{itemize}

\begin{tcolorbox}[breakable,colback=teal!4,colframe=teal!35!black,title={Before you move on}]
What does \ar{أثمر} mean? (\textquotedblleft{}To bear fruit\textquotedblright{}.) Say it once more.
\end{tcolorbox}
